% ── 다리 초안: mckinnon1983 → 우리 eq:sm-bare(원형) + eq:Veq(상호작용형) ──
% 삽입 위치: ch1_sec02a_part0.tex, McKinnon–Haering 문단(현행 156–161행) 직후,
%   verifybox(현행 162행) 앞.
% 앵커 문장: "...삽입 전극에의 적용 원형은 McKinnon--Haering \cite{mckinnon1983} 이다."(현행 157–158행)
% 컨테이너: srcbox. 우리 기호로 서술.
% ★확인 필요: mckinnon1983 = 단행본 장(Modern Aspects of Electrochemistry No.15, pp.235–304),
%   웹 원문 미확보 → 정확한 장내 식번호 미확인. 격자기체 삽입 등온선의 표준형은 교과서 확립
%   (Langmuir/Nernst-lattice + 평균장 상호작용) — 대응은 그 표준 결과 구조로 제시.

\begin{srcbox}
\textbf{다리 --- eq:sm-bare 와 eq:Veq 가 어느 삽입 등온선인가.} 식~\eqref{eq:sm-bare} 의 두-상태 원형과
\S\ref{sec:hys} 의 상호작용형 식~\eqref{eq:Veq} 는 McKinnon--Haering 이 삽입 전극에 세운 격자기체
등온선\cite{mckinnon1983} 의 원형$\cdot$확장에 각각 대응한다:
\begin{center}\footnotesize
\begin{tabular}{@{}ll@{}}
\hline
본문 & McKinnon--Haering\cite{mckinnon1983} 대응(격자기체) \\
\hline
평균 점유 $\langle n\rangle^{0}$ (식~\eqref{eq:sm-bare}) & 삽입 분율 $x$ (자리 점유율) \\
$\tfrac{RT}{sF}\ln\tfrac{\xi}{1-\xi}$ 항 (식~\eqref{eq:Veq}) & 배치 항 $\tfrac{kT}{e}\ln\tfrac{x}{1-x}$ \\
표준 중심 $U_j$ & 표준 전위 $E_0$ \\
상호작용 몫 $\tfrac{\Omega_j}{sF}(1-2\xi)$ (식~\eqref{eq:Veq}) & 평균장 자리-자리 상호작용 항 \\
문턱 $\Omega_j>2RT$ (식~\eqref{eq:spinodal}) & 상호작용이 1차 전이(전위 계단)를 낳는 임계 세기 \\
\hline
\end{tabular}
\end{center}
\emph{방법 요지} --- McKinnon--Haering 은 삽입 자리를 동등한 격자기체로 보고 각 자리 점유를 대정준(Langmuir
형) 평형으로 닫아 삽입 등온선(전위--조성 관계)을 얻고, 자리 간 평균장 상호작용을 더해 그 세기가 임계를
넘으면 1차 전이(전위 계단$\cdot$plateau)가 나타남을 보인다. \emph{가정 차} --- 원 원형은 점유 $0/1$ 만으로
미시상태를 지정하나(내부 자유도 없음), 본문은 점유 자리에 진동 등 내부 자유도 $q(T)$(식~\eqref{eq:partfn})를
실어 유효 자리에너지 $\tilde\varepsilon(T){=}\varepsilon_0{-}\kB T\ln q(T)$ 로 온도 이동을 허용하고(Part T 의
엔트로피$\cdot$발열 연결), 상호작용도 규칙용액 $\Omega_j$ 로 파라미터화해 spinodal 문턱 $\Omega_j{>}2RT$ 를
명시 분기로 닫는다.
\end{srcbox}
